\subsubsection{Floating-Point Exception Causes and Commit}

\begin{BedrockTabular}{@{}>{\raggedright\arraybackslash}p{0.55in}>{\raggedright\arraybackslash}p{4.85in}@{}}
\toprule
\rowcolor{ManualHeaderFill}
\textbf{Cause} & \textbf{Generated when}\\
\midrule
\texttt{NV} & an input is a signaling NaN; an arithmetic operation has no defined numeric result, including infinity minus the same infinity, zero times infinity, zero divided by zero, infinity divided by infinity, or square root of a negative nonzero value; or a floating-to-integer conversion is invalid\\
\texttt{DZ} & a finite nonzero value is divided by zero\\
\texttt{OF} & the rounded numeric magnitude exceeds the selected format\textquotesingle{}s finite range\\
\texttt{UF} & the result is tiny after rounding and inexact, or FTZ flushes a nonzero subnormal result\\
\texttt{NX} & rounding or a valid conversion changes the exact value, or OF or FTZ generates NX\\
\bottomrule
\end{BedrockTabular}

If any generated floating-point exception cause has its matching FSTATUS enable bit set, the instruction raises the
(:ref:FP.events.EXCEPTION.FLOATING_POINT_EXCEPTION:) architectural event. Its error-code bitmap contains every cause generated by the operation.
The destination, integer FLAGS, and FFLAGS remain unchanged, and arithmetic instructions never modify FSTATUS.

If no generated floating-point exception cause is enabled, the instruction commits its destination and any complete
integer FLAGS image defined by the instruction, and ORs every generated cause into FFLAGS. FFLAGS (:term:base.terms.field|plural:) not named by
the instruction remain unchanged. An instruction-specific rule may exclude a cause.
